\section{The released result and its actual signed path}
\label{nsprop:source}
The source in this section is the supplied 166-page edition of
\emph{Finite Time Blowup for Navier--Stokes}, whose SHA-256 is
\begin{center}\footnotesize
\texttt{0e779481c4da40bd28d1e642e1d8ca57447d129610df28dfa5a11e9af8ae228f}.
\end{center}
Its Theorem 1.1, on page 1, states the following result for the standard
three-dimensional incompressible Navier--Stokes equation. For every
fixed physical viscosity $\nu_{\rm NS}>0$, the source constructs
\[
 f\in C^\infty_c(\mathbb R^3\times(0,\infty);\mathbb R^3),
 \qquad u,p\in C^\infty(\mathbb R^3\times[0,1)),
\]
with a fixed compact spatial support $K$ for $u$ and $p$, such that
\begin{equation}
 \begin{gathered}
 \partial_tu+(u\cdot\nabla)u-\nu_{\rm NS}\Delta u+\nabla p=f,
 \quad \nabla\cdot u=0,\quad u(\cdot,0)=0,\\
 \quad \sup_{t<1}\|u(t)\|_2<\infty,\quad
 \limsup_{t\uparrow1}\|u(t)\|_\infty=\infty .
 \end{gathered}
 \label{nsprop:source_theorem}
\end{equation}
It also excludes a global smooth solution with the same force and datum
and uniformly bounded kinetic energy. Its Corollary 10.6, pages 125--126,
gives a periodic construction on $\mathbb R^3/\mathbb Z^3$ with smooth
time-compact force, zero initial datum and no global smooth continuation.
The source identifies these with alternatives (C) and (D) in the
referenced Millennium problem statement. The force is nonzero, as the
source explicitly notes on page 124. These are statements of the cited
source theorem; the finite algebra checks in this lane do not constitute
an independent verification of the complete 166-page proof.

The source proves more than an unsigned norm divergence. In its
viscosity-one construction, equations (10.20)--(10.21), page 124, give
\begin{equation}
 \begin{split}
 A&=\tfrac12+h,\qquad 0<h<1/100,\qquad X_{\rm in}>0,\quad e_0>0,\\
 x_\tau&=(\sqrt{2X_{\rm in}\tau},0,0),\qquad t_\tau=1-\tau,\\
 u_\theta(x_\tau,t_\tau)
   &=\tau^{-A}(e_0+\epsilon(\tau)),\qquad
      |\epsilon(\tau)|\le C\tau^{2h}\quad(0<\tau<\tau_0).
 \end{split}\label{nsprop:positive_path}
\end{equation}
Here $C,\tau_0$ are finite constants supplied by the source asymptotic;
$e_0$ and $X_{\rm in}$ retain their source values. If $C>0$, choose
$0<\tau_1\le\tau_0$ so that $C\tau_1^{2h}\le e_0/2$; if $C=0$,
take $\tau_1=\tau_0$. It follows directly that
\[
 u_\theta(x_\tau,t_\tau)\ge \frac{e_0}{2}\tau^{-A}
       \longrightarrow+\infty\qquad(0<\tau<\tau_1).
\]
The observation point moves to the origin. This assertion does not
replace that path by the fixed point $x=0$. On the path, $q=\tau>0$,
and at its cylindrical angle $\theta=0$ one has $e_\theta=e_2$.
The geometric domain $q>0$, the positive exponent $h$, and the signed
velocity value are three retained parts of the same explicit formula.

\section{Exact propagation at retained viscosities and orientations}
\label{nsprop:map}
For any smooth fields $u,p$ define the physical momentum residual
\[
 R_{\nu_b}(u,p)=\partial_tu+(u\cdot\nabla)u-\nu_b\Delta u+\nabla p.
\]
Let $\nu_b>0$ and $\nu_t>0$ be the base and target physical viscosities,
and retain
\[
 \rho=\sqrt{\nu_t/\nu_b},\qquad O^TO=OO^T=I,\qquad a\in\mathbb R^3.
\]
Define the affine spatial coordinate $y=O^T(x-a)/\rho$ and the fields
\begin{equation}
 u^\sharp(x,t)=\rho O u(y,t),\quad
 p^\sharp(x,t)=\rho^2p(y,t),\quad
 f^\sharp(x,t)=\rho O f(y,t).
 \label{nsprop:affine_map}
\end{equation}
No time coordinate is changed. These formulas give a bijection on smooth
field triples, with inverse
\[
 u(y,t)=\rho^{-1}O^Tu^\sharp(a+\rho Oy,t),\quad
 p(y,t)=\rho^{-2}p^\sharp(a+\rho Oy,t),\quad
 f(y,t)=\rho^{-1}O^Tf^\sharp(a+\rho Oy,t).
\]
For the Jacobian convention $(Du)_{ij}=\partial_j u_i$, differentiation
gives, without dropping an orientation or viscosity,
\begin{align*}
 Du^\sharp&=O(Du)(y,t)O^T,&
 \nabla\cdot u^\sharp&=(\nabla\cdot u)(y,t),\\
 \partial_tu^\sharp&=\rho O\partial_tu(y,t),&
 (u^\sharp\cdot\nabla_x)u^\sharp&=\rho O((u\cdot\nabla_y)u)(y,t),\\
 \Delta_xu^\sharp&=\rho^{-1}O\Delta_yu(y,t),&
 \nabla_xp^\sharp&=\rho O\nabla_yp(y,t).
\end{align*}
The nonlinear identity follows by multiplying $Du^\sharp$ by
$u^\sharp$ and using $O^TO=I$. The Laplacian identity follows by
contracting the two spatial derivative indices and using orthogonality.
Since $\nu_t=\rho^2\nu_b$, these identities prove
\begin{equation}
 R_{\nu_t}(u^\sharp,p^\sharp)-f^\sharp
       =\rho O\bigl(R_{\nu_b}(u,p)-f\bigr)(y,t).
 \label{nsprop:residual_map}
\end{equation}
In particular the map preserves the exact equation as well as each
nonzero residual with its exact factor.

For vorticity, the cross-product identity
$(Ov)\times(Ow)=\det(O)O(v\times w)$ gives
\begin{equation}
 \omega^\sharp(x,t)=\det(O)O\omega(y,t).
 \label{nsprop:vorticity}
\end{equation}
One proof of this identity takes its scalar product with $Oz$:
both sides give $\det(O)\det[v,w,z]$ for every $z$. Applying it to
the contracted first derivatives proves the vorticity formula; the
factor $\rho$ in velocity cancels the inverse factor in a derivative.

Spatial support becomes $a+\rho OK$. The Jacobian has absolute
determinant $\rho^3$, for both signs of $\det O$. Hence, for
$1\le p<\infty$,
\begin{equation}
 \|u^\sharp(t)\|_p=\rho^{1+3/p}\|u(t)\|_p,\qquad
 \|u^\sharp(t)\|_\infty=\rho\|u(t)\|_\infty,\qquad
 \|u^\sharp(t)\|_2^2=\rho^5\|u(t)\|_2^2 .
 \label{nsprop:norms}
\end{equation}
The Frobenius norm of $O(Du)O^T$ equals that of $Du$. Consequently,
with integrals over the same time interval,
\[
 \nu_t\int\|\nabla u^\sharp\|_2^2\,dt
     =\rho^5\nu_b\int\|\nabla u\|_2^2\,dt .
\]
For arbitrary directions $v_1,\ldots,v_j$, the exact force derivative is
\[
 D_x^j\partial_t^m f^\sharp[v_1,\ldots,v_j](x,t)
 =\rho^{1-j}O
   (D_y^j\partial_t^m f)[O^Tv_1,\ldots,O^Tv_j](y,t).
\]
Thus smoothness, compact space-time support and vanishing of every
derivative at a terminal point are preserved. The initial datum remains
zero. A global smooth bounded-energy competitor for the transformed
datum and force would, under the displayed inverse, be a competitor
for the base datum and force, because its energy is multiplied by the
fixed factor $\rho^{-5}$. This proves propagation of the source
noncontinuation conclusion through the map.

\section{The source positive branch and the exact negative branch}
\label{nsprop:branches}
Apply the map to the source field in \eqref{nsprop:positive_path}.
Here $\nu_b=1$ specifies the source object being used; physical target
viscosity $\nu_t$ remains arbitrary and $\rho=\sqrt{\nu_t}$ is retained.
Take $a=0$ and the two particular orthogonal matrices
\[
 O_+=I,\qquad O_-=\operatorname{diag}(1,-1,1).
\]
Both fix $e_1$, so they give the same spatial observation path
$x^\sharp_\tau=\rho x_\tau$ and the same time $1-\tau$.
Since $O_+e_2=e_2$ and $O_-e_2=-e_2$, their exact signed values are
\begin{equation}
 u^\sharp_{\pm,\theta}(\rho x_\tau,1-\tau)
     =\pm\rho\,\tau^{-A}(e_0+\epsilon(\tau))
        \longrightarrow\begin{cases}+\infty,&+,\\-\infty,&-.\end{cases}
 \label{nsprop:two_paths}
\end{equation}
All support, energy, force and viscosity statements in the preceding
section apply to both triples. For the negative branch the force is
$\rho O_-f(O_-x/\rho,t)$. Equality of this force with the positive
branch force is not needed and is not asserted.

For a general, possibly nonaxisymmetric field the reflection alone has
the exact cylindrical component formula
\[
 (u^r,u^\theta,u^z)^\sharp(r,\theta,z,t)
    =(u^r,-u^\theta,u^z)(r,-\theta,z,t).
\]
Indeed $O_-e_r(-\theta)=e_r(\theta)$,
$O_-e_\theta(-\theta)=-e_\theta(\theta)$ and $O_-e_z=e_z$.
This proves the formula even in the presence of all oscillatory
corrections. For the axisymmetric background the argument $-\theta$
can be dropped because its components are independent of $\theta$.
For a velocity covariance tensor the exact Cartesian transformation is
$T^\sharp(x,t)=\rho^2 O T(y,t)O^T$: expand the outer product of
$\rho Ow$ with itself and use linearity of the retained average.
For its angularly averaged symmetric tensor, the reflection basis
calculation at $\rho=1$ gives
\[
 (T_{r\theta},T_{rz})^\sharp=(-T_{r\theta},T_{rz}).
\]
Under the additional viscosity dilation its components have factor
$\rho^2$ at the corresponding points. Thus the two covariance columns
and their target transform together. Positive coefficients in a source
representation $T=c_1v_1+c_2v_2$, $c_i>0$, remain the same coefficients
of the transformed columns; no change of coefficient sign is required
to realize the negative-swirl branch.

The source's periodic construction also has this reflection: $O_-$
preserves $\mathbb Z^3$, so $[x]\mapsto[O_-x]$ is well defined on
$\mathbb R^3/\mathbb Z^3$. The derivative identities apply to periodic
lifts. This gives both signed orientations of the source Corollary 10.6.
Arbitrary $O\in O(3)$ need not preserve this particular lattice; only
the lattice-preserving reflection just specified is used for this
periodic assertion.

\section{Propagation into the coupled calculation}
\label{nsprop:programme}
The preceding sections use the released whole-space and periodic
constructions at the precise places where those source results apply.
The original coupled Boussinesq work is due to Levent Alp\"oge and
Tristan Buckmaster. This lane's retained coupled system, its finite
controlled stages and its residual maps remain the separate objects
defined in the preceding parts of this reader.

For every compact candidate $(u_*+w_y,p_*+\pi_y)$ in the covariance
part, \eqref{nsprop:residual_map} propagates its full computed residual.
Linearity of the transformation and the exact nonlinear identity give
in particular
\[
 \Delta R_y^\sharp(x,t)=\rho O\Delta R_y(y,t).
\]
Thus the finite candidate's cutoff forces, curl errors and covariance
remainders all have an exact transformed counterpart. A nonzero
remainder is carried to a nonzero remainder, because the map is
invertible. The source's residual flatness at the singular point belongs
to its completed local source construction; its extended smooth force
need not vanish on the entire time-one slice. The next section calculates the
surviving moments of the actual finite candidate, so that propagation
continues through its computed terms rather than assuming that its
different coupled coefficients inherit the source's infinite estimates.

The earlier suggestion of a negative blowout imposes no condition on
the programme. Formula \eqref{nsprop:positive_path} is the source-selected
sign; formula \eqref{nsprop:two_paths} proves the exact relation of its
two orientations. Both formulas retain a positive physical viscosity
and the source coordinate domain $q>0$.
